% context: markscheme (official solutions) for 1-5CW_Applications
% topic: unit conversion chains, scientific notation, significant figures,
%        estimation
% convention: "=" joins exact equalities; "\approx" introduces a rounded value
\documentclass[12pt, twoside]{article}
\input{../preamble}
\usepackage{float}
\setlength{\parskip}{0.3\baselineskip}

\title{Physics}
\author{Chris Huson}
\date{September 2026}

\fancyhead[LE]{\thepage}
\fancyhead[RO]{\thepage}
\fancyhead[LO]{La Scuola d'Italia / Huson / Physics \\* 25 September 2026}

\begin{document}

\subsubsection*{1.5 Classwork Solutions: Applications}

\emph{Notation: \,$=$\, joins values that are exactly equal; \,$\approx$\, introduces a
rounded value. Each conversion runs as one chain of unit factors and is evaluated once,
so there is nothing to round in the middle.}

\begin{enumerate}[itemsep=0.3cm]

\item Rain on Manhattan.
  \begin{enumerate}[itemsep=0.15cm]
    \item $V = 20{,}000~\text{m} \times 3000~\text{m} \times 0.025~\text{m}
      = \mathbf{1.5 \times 10^{6}~\textbf{m}^3}$

      \emph{Converting each length to meters first avoids the square-kilometer trap:
      $1~\text{km}^2$ is $10^{6}~\text{m}^2$, not $1000~\text{m}^2$.}
    \item $1.5 \times 10^{6}~\text{m}^3 \times \dfrac{1000~\text{L}}{1~\text{m}^3}
      \times \dfrac{1~\text{kg}}{1~\text{L}} = \mathbf{1.5 \times 10^{9}~\textbf{kg}}$

      $1.5 \times 10^{9}~\text{kg} \times \dfrac{1~\text{t}}{1000~\text{kg}}
      = \mathbf{1.5 \times 10^{6}~\textbf{metric tons}}$
    \item $1.5 \times 10^{6}~\text{m}^3 \times \dfrac{1000~\text{L}}{1~\text{m}^3}
      \times \dfrac{1~\text{gal}}{3.785~\text{L}} = 3.963\ldots \times 10^{8}~\text{gal}
      \approx \mathbf{4.0 \times 10^{8}~\textbf{gal}}$

      About \textbf{0.4 of a day's supply}, or roughly 10 hours of the whole city's
      water use --- from one storm, on one borough.

      \emph{Two figures throughout, from the 2.5 cm and 3.0 km (the 20 km is a rough
      rectangle too). Start the gallons from the volume in part (a), not from a rounded
      intermediate.}
  \end{enumerate}

\item The oil slick.
  \begin{enumerate}[itemsep=0.15cm]
    \item $1.0~\text{L} \times \dfrac{1~\text{m}^3}{1000~\text{L}}
      = \mathbf{1.0 \times 10^{-3}~\textbf{m}^3}$
    \item $A = \dfrac{V}{\text{thickness}}
      = \dfrac{1.0 \times 10^{-3}~\text{m}^3}{2 \times 10^{-10}~\text{m}}
      = \mathbf{5 \times 10^{6}~\textbf{m}^2}$

      \emph{Dividing the powers of ten: $10^{-3-(-10)} = 10^{7}$, and $1.0/2 = 0.5$. The
      units work too: $\text{m}^3 / \text{m} = \text{m}^2$.}
    \item $\pi r^2 = 5 \times 10^{6}~\text{m}^2$, so
      $r = \sqrt{\dfrac{5 \times 10^{6}~\text{m}^2}{\pi}} = 1261.5\ldots~\text{m}$
      and $d = 2r = 2523\ldots~\text{m} \approx$ \textbf{2.5 km} (about 3 km)

      \emph{The molecule size has one significant figure, so this is an estimate: ``a
      couple of kilometers'' is the physics. Franklin's teaspoon covered about half an
      acre. In 1890 Lord Rayleigh repeated the experiment carefully and used it to make
      one of the first estimates of the size of a molecule.}
  \end{enumerate}

\item Walking around the Earth.
  \begin{enumerate}[itemsep=0.15cm]
    \item $40{,}075~\text{km} \times \dfrac{1~\text{h}}{5.0~\text{km}}
      \times \dfrac{1~\text{day}}{12~\text{h}} = 667.91\ldots~\text{days}
      \approx \mathbf{670~\textbf{days}}$, better written
      $\mathbf{6.7 \times 10^{2}~\textbf{days}}$

      \emph{Two figures, set by the 5.0 km/h. The 12 hours is a stated assumption,
      not a measurement.}
    \item $667.91~\text{days} \times \dfrac{1~\text{yr}}{365~\text{days}}
      = 1.829\ldots~\text{yr} \approx \mathbf{1.8~\textbf{years}}$

      \emph{Carry the unrounded 667.91 into this step. And the Pacific Ocean is in the
      way.}
  \end{enumerate}

\item Sunlight.

  $t = 1.496 \times 10^{8}~\text{km} \times \dfrac{1000~\text{m}}{1~\text{km}}
  \times \dfrac{1~\text{s}}{3.00 \times 10^{8}~\text{m}}
  \times \dfrac{1~\text{min}}{60~\text{s}}
  = 8.3111\ldots~\text{min} \approx \mathbf{8.31~\textbf{min}}$

  \emph{Three figures, set by the 3.00. Rearranging $v = d/t$ to $t = d/v$ is the same
  thing as the $\frac{1~\text{s}}{3.00 \times 10^{8}~\text{m}}$ factor --- the speed of
  light is a conversion factor between meters and seconds. The Sun we see is the Sun
  of eight minutes ago.}

\end{enumerate}
\end{document}
